% =============================================================================
% Assignments/CS401/05-addressing-cpu-bus-solutions.tex
% Assignment 5 Solution Key: Addressing Modes, CPU Organization, and Bus Structures
%
% Mirrors 05-addressing-cpu-bus-assignment.tex's numbering. Instructor/answer-
% key file — do not distribute to students alongside the assignment.
%
% Compile with: cd Assignments/CS401 &&
%   TEXINPUTS=../../Preambles:$TEXINPUTS xelatex -interaction=nonstopmode 05-addressing-cpu-bus-solutions.tex
%   BIBINPUTS=../..:$BIBINPUTS bibtex 05-addressing-cpu-bus-solutions
%   TEXINPUTS=../../Preambles:$TEXINPUTS xelatex -interaction=nonstopmode 05-addressing-cpu-bus-solutions.tex (x2)
% =============================================================================

\documentclass[11pt]{article}
\usepackage[margin=1in]{geometry}
% =============================================================================
% Preambles/header.tex — shared LaTeX/Beamer preamble for this project
%
% Use from a Beamer lecture by prepending:
%     % =============================================================================
% Preambles/header.tex — shared LaTeX/Beamer preamble for this project
%
% Use from a Beamer lecture by prepending:
%     % =============================================================================
% Preambles/header.tex — shared LaTeX/Beamer preamble for this project
%
% Use from a Beamer lecture by prepending:
%     \input{header}
% and compiling with TEXINPUTS=../Preambles:$TEXINPUTS (the /compile-latex
% skill does this automatically).
%
% PALETTE CONTRACT. Color names defined here MUST match the names in
% Quarto/theme-template.scss so Beamer and Quarto renderings use the same
% palette. The scripts/check-palette-sync.sh script enforces this.
%
% To customize: edit the HEX values below AND the matching $variable in
% Quarto/theme-template.scss, then run scripts/check-palette-sync.sh.
%
% TITLE PAGE. A formal, serif (Libertine) title-page template is defined in
% the Beamer-only block below (\setbeamertemplate{title page}) -- title,
% subtitle, a thin gold rule, then \author/\institute/\date. Frametitle and
% body text remain on Beamer's sans-serif default. No SCSS counterpart is
% needed for this (font/layout only, no new colors).
% =============================================================================

% -----------------------------------------------------------------------------
% Engine detection + encoding
%
% Our /compile-latex skill uses XeLaTeX, where inputenc is unnecessary and
% can produce encoding warnings. Only load inputenc under pdfTeX.
% -----------------------------------------------------------------------------
\usepackage{iftex}
\ifPDFTeX
  \usepackage[utf8]{inputenc}
\fi

\usepackage{amsmath, amssymb, amsthm}
\usepackage{booktabs}
\usepackage{graphicx}

% -----------------------------------------------------------------------------
% Serif face for the title page (formal/academic look). Linux Libertine is
% XeLaTeX- and pdfTeX-safe and redefines \rmdefault, so \rmfamily picks it up
% everywhere \rmfamily is used explicitly. Body text and frametitles stay on
% Beamer's own sans-serif default (unaffected) -- only the title-page
% \setbeamerfont{...} overrides below opt in to \rmfamily.
% -----------------------------------------------------------------------------
\usepackage{libertine}

% -----------------------------------------------------------------------------
% PALETTE — mirrors Quarto/theme-template.scss variable names.
% Load xcolor and define colors BEFORE anything that references them
% (notably hyperref, hypersetup, and the Beamer color assignments).
% -----------------------------------------------------------------------------
\usepackage{xcolor}

\definecolor{primary-blue}{HTML}{012169}      % headings, accents
\definecolor{primary-gold}{HTML}{B9975B}      % emphasis, borders
\definecolor{highlight-yellow}{HTML}{F2A900}  % markers, alerts
\definecolor{light-bg}{HTML}{E8EDF5}          % title / section backgrounds
\definecolor{jet}{HTML}{1A1A1A}               % body text

% Semantic colors (used in diagrams and annotations)
\definecolor{positive}{HTML}{15803D}          % good / observed / identified
\definecolor{negative}{HTML}{B91C1C}          % bad / problematic / bias
\definecolor{neutral}{HTML}{525252}           % reference / context

% Highlight accents (match .hi-* classes in the SCSS)
\definecolor{hi-slate}{HTML}{314F4F}
\definecolor{hi-green}{HTML}{2E8B57}
\definecolor{hi-red}{HTML}{C41E3A}

% Semantic aliases so skill templates and snippets can write
% \textcolor{accent}{...} without hard-coding "primary-blue".
\colorlet{accent}{primary-blue}
\colorlet{accent2}{primary-gold}

% -----------------------------------------------------------------------------
% Hyperref — loaded AFTER xcolor + palette so link colors resolve.
% -----------------------------------------------------------------------------
\usepackage{hyperref}
\hypersetup{colorlinks=true, linkcolor=primary-blue, urlcolor=primary-blue,
            citecolor=primary-blue, pdfborder={0 0 0}}

% -----------------------------------------------------------------------------
% Course code — set once per deck (after \input{header}, before \begin{document})
% via \coursecode{CS401}. Shown in the footer on every slide so a deck's course
% is identifiable without opening the title page. Empty by default.
% -----------------------------------------------------------------------------
\makeatletter
\newcommand{\coursecode}[1]{\gdef\@coursecodetext{#1}}
\gdef\@coursecodetext{}
\makeatother

% -----------------------------------------------------------------------------
% Beamer theme assignments (only apply under Beamer).
%
% We detect Beamer via \@ifclassloaded{beamer}, which is the documented,
% non-internal way. This must be wrapped in \makeatletter/\makeatother
% because \@ifclassloaded contains an @-catcode character.
% -----------------------------------------------------------------------------
\makeatletter
\@ifclassloaded{beamer}{%
  \usefonttheme{professionalfonts}
  \setbeamercolor{structure}{fg=primary-blue}
  \setbeamercolor{frametitle}{fg=primary-blue}
  \setbeamercolor{title}{fg=primary-blue}
  \setbeamercolor{subtitle}{fg=primary-gold}
  \setbeamercolor{author}{fg=primary-blue}
  \setbeamercolor{institute}{fg=neutral}
  \setbeamercolor{date}{fg=primary-gold}
  \setbeamercolor{itemize item}{fg=primary-blue}
  \setbeamercolor{itemize subitem}{fg=primary-gold}
  \setbeamercolor{alerted text}{fg=primary-gold}
  \setbeamercolor{block title}{fg=white, bg=primary-blue}
  \setbeamercolor{block body}{bg=light-bg}
  % Semantic block variants: block=definition/concept (blue), exampleblock=
  % worked example (green/positive), alertblock=key takeaway or caution (gold).
  % Keep to <=2 colored boxes per slide across ALL three variants (INV-7).
  \setbeamercolor{block title example}{fg=white, bg=positive}
  \setbeamercolor{block body example}{bg=light-bg}
  \setbeamercolor{block title alerted}{fg=white, bg=primary-gold}
  \setbeamercolor{block body alerted}{bg=light-bg}

  % ---------------------------------------------------------------------------
  % Formal title page: serif (Libertine, via \rmfamily) for title-page
  % elements only. Body text, itemize, and frametitle are intentionally left
  % on Beamer's sans-serif default -- \usebeamerfont{title} is also what
  % \transitionslide (below) uses for its big centered text, so section
  % transitions pick up the same serif treatment for free.
  % ---------------------------------------------------------------------------
  \setbeamerfont{title}{family=\rmfamily, series=\bfseries, size=\Large}
  \setbeamerfont{subtitle}{family=\rmfamily, shape=\itshape, size=\normalsize}
  \setbeamerfont{author}{family=\rmfamily, size=\normalsize}
  \setbeamerfont{institute}{family=\rmfamily, size=\scriptsize}
  \setbeamerfont{date}{family=\rmfamily, size=\scriptsize}

  \setbeamertemplate{title page}{
    \vfill
    \begin{center}
      {\usebeamerfont{title}\usebeamercolor[fg]{title}\inserttitle\par}
      \vspace{0.15cm}
      {\usebeamerfont{subtitle}\usebeamercolor[fg]{subtitle}\insertsubtitle\par}
      \vspace{0.45cm}
      {\color{primary-gold}\rule{0.42\textwidth}{0.6pt}\par}
      \vspace{0.45cm}
      {\usebeamerfont{author}\usebeamercolor[fg]{author}\insertauthor\par}
      \vspace{0.35cm}
      {\usebeamerfont{institute}\usebeamercolor[fg]{institute}\insertinstitute\par}
      \vspace{0.35cm}
      {\usebeamerfont{date}\usebeamercolor[fg]{date}\insertdate\par}
    \end{center}
    \vfill
  }

  % Minimal footer: course code (left) + page X / N (right), no navigation clutter
  \setbeamertemplate{navigation symbols}{}
  \setbeamertemplate{footline}{%
    \color{neutral}\footnotesize\hspace{0.7em}\@coursecodetext\hfill
    \insertframenumber/\inserttotalframenumber\hspace{0.7em}\vspace{0.3em}}
}{}
\makeatother

% -----------------------------------------------------------------------------
% TikZ setup — libraries the snippet gallery uses
% -----------------------------------------------------------------------------
\usepackage{tikz}
\usetikzlibrary{arrows.meta, positioning, calc, decorations.pathreplacing,
                fit, shapes.geometric, backgrounds}

% Shared TikZ styles — snippets and hand-written diagrams can pick these up
% via `\begin{tikzpicture}[every node/.style={...}]` or by naming specific
% styles (e.g., `\node[dag-node] (X) at (0,0) {X};`).
\tikzset{
  % Node styles for causal DAGs and similar
  dag-node/.style = {
    draw=primary-blue, fill=light-bg, rounded corners=2pt,
    minimum width=1.2cm, minimum height=0.9cm, align=center, font=\small
  },
  decision-node/.style = {
    draw=primary-gold, fill=white, diamond, aspect=1.8,
    minimum width=1.8cm, minimum height=1.2cm, align=center, font=\small,
    inner sep=1pt
  },
  % Edge styles — observed vs counterfactual
  observed-edge/.style = {-{Latex[length=2.5mm]}, thick, draw=jet},
  counterfactual-edge/.style = {-{Latex[length=2.5mm]}, thick, draw=neutral, dashed},
  confound-edge/.style = {-{Latex[length=2.5mm]}, thick, draw=negative, dashed},
  % Data-point styles for plots
  observed-dot/.style = {fill=primary-blue, draw=primary-blue, circle, inner sep=1.2pt},
  counterfactual-dot/.style = {fill=white, draw=neutral, circle, inner sep=1.2pt}
}

% -----------------------------------------------------------------------------
% Convenience macros
% -----------------------------------------------------------------------------
\newcommand{\muted}[1]{\textcolor{neutral}{#1}}
\newcommand{\key}[1]{\textcolor{primary-gold}{\textbf{#1}}}
\newcommand{\good}[1]{\textcolor{positive}{#1}}
\newcommand{\bad}[1]{\textcolor{negative}{#1}}

% Standout frame macro: full-bleed dark background for section transitions.
% Beamer-only (uses \begin{frame}/\setbeamercolor/\usebeamerfont) -- do not
% call from an article-class file (e.g. Notes/<CODE>/*-notes.tex). \muted,
% \key, \good, \bad, and \coursecode above are plain xcolor/gdef and are
% safe to use from either class; verified when Notes/article-class support
% was added.
% Expand: \transitionslide{Your transition title}
\newcommand{\transitionslide}[1]{%
  \begingroup
  \setbeamercolor{background canvas}{bg=primary-blue}
  \begin{frame}[plain]
    \centering
    \vfill
    {\usebeamerfont{title}\color{white}\Large #1\par}
    \vfill
  \end{frame}
  \endgroup
}

% and compiling with TEXINPUTS=../Preambles:$TEXINPUTS (the /compile-latex
% skill does this automatically).
%
% PALETTE CONTRACT. Color names defined here MUST match the names in
% Quarto/theme-template.scss so Beamer and Quarto renderings use the same
% palette. The scripts/check-palette-sync.sh script enforces this.
%
% To customize: edit the HEX values below AND the matching $variable in
% Quarto/theme-template.scss, then run scripts/check-palette-sync.sh.
%
% TITLE PAGE. A formal, serif (Libertine) title-page template is defined in
% the Beamer-only block below (\setbeamertemplate{title page}) -- title,
% subtitle, a thin gold rule, then \author/\institute/\date. Frametitle and
% body text remain on Beamer's sans-serif default. No SCSS counterpart is
% needed for this (font/layout only, no new colors).
% =============================================================================

% -----------------------------------------------------------------------------
% Engine detection + encoding
%
% Our /compile-latex skill uses XeLaTeX, where inputenc is unnecessary and
% can produce encoding warnings. Only load inputenc under pdfTeX.
% -----------------------------------------------------------------------------
\usepackage{iftex}
\ifPDFTeX
  \usepackage[utf8]{inputenc}
\fi

\usepackage{amsmath, amssymb, amsthm}
\usepackage{booktabs}
\usepackage{graphicx}

% -----------------------------------------------------------------------------
% Serif face for the title page (formal/academic look). Linux Libertine is
% XeLaTeX- and pdfTeX-safe and redefines \rmdefault, so \rmfamily picks it up
% everywhere \rmfamily is used explicitly. Body text and frametitles stay on
% Beamer's own sans-serif default (unaffected) -- only the title-page
% \setbeamerfont{...} overrides below opt in to \rmfamily.
% -----------------------------------------------------------------------------
\usepackage{libertine}

% -----------------------------------------------------------------------------
% PALETTE — mirrors Quarto/theme-template.scss variable names.
% Load xcolor and define colors BEFORE anything that references them
% (notably hyperref, hypersetup, and the Beamer color assignments).
% -----------------------------------------------------------------------------
\usepackage{xcolor}

\definecolor{primary-blue}{HTML}{012169}      % headings, accents
\definecolor{primary-gold}{HTML}{B9975B}      % emphasis, borders
\definecolor{highlight-yellow}{HTML}{F2A900}  % markers, alerts
\definecolor{light-bg}{HTML}{E8EDF5}          % title / section backgrounds
\definecolor{jet}{HTML}{1A1A1A}               % body text

% Semantic colors (used in diagrams and annotations)
\definecolor{positive}{HTML}{15803D}          % good / observed / identified
\definecolor{negative}{HTML}{B91C1C}          % bad / problematic / bias
\definecolor{neutral}{HTML}{525252}           % reference / context

% Highlight accents (match .hi-* classes in the SCSS)
\definecolor{hi-slate}{HTML}{314F4F}
\definecolor{hi-green}{HTML}{2E8B57}
\definecolor{hi-red}{HTML}{C41E3A}

% Semantic aliases so skill templates and snippets can write
% \textcolor{accent}{...} without hard-coding "primary-blue".
\colorlet{accent}{primary-blue}
\colorlet{accent2}{primary-gold}

% -----------------------------------------------------------------------------
% Hyperref — loaded AFTER xcolor + palette so link colors resolve.
% -----------------------------------------------------------------------------
\usepackage{hyperref}
\hypersetup{colorlinks=true, linkcolor=primary-blue, urlcolor=primary-blue,
            citecolor=primary-blue, pdfborder={0 0 0}}

% -----------------------------------------------------------------------------
% Course code — set once per deck (after % =============================================================================
% Preambles/header.tex — shared LaTeX/Beamer preamble for this project
%
% Use from a Beamer lecture by prepending:
%     \input{header}
% and compiling with TEXINPUTS=../Preambles:$TEXINPUTS (the /compile-latex
% skill does this automatically).
%
% PALETTE CONTRACT. Color names defined here MUST match the names in
% Quarto/theme-template.scss so Beamer and Quarto renderings use the same
% palette. The scripts/check-palette-sync.sh script enforces this.
%
% To customize: edit the HEX values below AND the matching $variable in
% Quarto/theme-template.scss, then run scripts/check-palette-sync.sh.
%
% TITLE PAGE. A formal, serif (Libertine) title-page template is defined in
% the Beamer-only block below (\setbeamertemplate{title page}) -- title,
% subtitle, a thin gold rule, then \author/\institute/\date. Frametitle and
% body text remain on Beamer's sans-serif default. No SCSS counterpart is
% needed for this (font/layout only, no new colors).
% =============================================================================

% -----------------------------------------------------------------------------
% Engine detection + encoding
%
% Our /compile-latex skill uses XeLaTeX, where inputenc is unnecessary and
% can produce encoding warnings. Only load inputenc under pdfTeX.
% -----------------------------------------------------------------------------
\usepackage{iftex}
\ifPDFTeX
  \usepackage[utf8]{inputenc}
\fi

\usepackage{amsmath, amssymb, amsthm}
\usepackage{booktabs}
\usepackage{graphicx}

% -----------------------------------------------------------------------------
% Serif face for the title page (formal/academic look). Linux Libertine is
% XeLaTeX- and pdfTeX-safe and redefines \rmdefault, so \rmfamily picks it up
% everywhere \rmfamily is used explicitly. Body text and frametitles stay on
% Beamer's own sans-serif default (unaffected) -- only the title-page
% \setbeamerfont{...} overrides below opt in to \rmfamily.
% -----------------------------------------------------------------------------
\usepackage{libertine}

% -----------------------------------------------------------------------------
% PALETTE — mirrors Quarto/theme-template.scss variable names.
% Load xcolor and define colors BEFORE anything that references them
% (notably hyperref, hypersetup, and the Beamer color assignments).
% -----------------------------------------------------------------------------
\usepackage{xcolor}

\definecolor{primary-blue}{HTML}{012169}      % headings, accents
\definecolor{primary-gold}{HTML}{B9975B}      % emphasis, borders
\definecolor{highlight-yellow}{HTML}{F2A900}  % markers, alerts
\definecolor{light-bg}{HTML}{E8EDF5}          % title / section backgrounds
\definecolor{jet}{HTML}{1A1A1A}               % body text

% Semantic colors (used in diagrams and annotations)
\definecolor{positive}{HTML}{15803D}          % good / observed / identified
\definecolor{negative}{HTML}{B91C1C}          % bad / problematic / bias
\definecolor{neutral}{HTML}{525252}           % reference / context

% Highlight accents (match .hi-* classes in the SCSS)
\definecolor{hi-slate}{HTML}{314F4F}
\definecolor{hi-green}{HTML}{2E8B57}
\definecolor{hi-red}{HTML}{C41E3A}

% Semantic aliases so skill templates and snippets can write
% \textcolor{accent}{...} without hard-coding "primary-blue".
\colorlet{accent}{primary-blue}
\colorlet{accent2}{primary-gold}

% -----------------------------------------------------------------------------
% Hyperref — loaded AFTER xcolor + palette so link colors resolve.
% -----------------------------------------------------------------------------
\usepackage{hyperref}
\hypersetup{colorlinks=true, linkcolor=primary-blue, urlcolor=primary-blue,
            citecolor=primary-blue, pdfborder={0 0 0}}

% -----------------------------------------------------------------------------
% Course code — set once per deck (after \input{header}, before \begin{document})
% via \coursecode{CS401}. Shown in the footer on every slide so a deck's course
% is identifiable without opening the title page. Empty by default.
% -----------------------------------------------------------------------------
\makeatletter
\newcommand{\coursecode}[1]{\gdef\@coursecodetext{#1}}
\gdef\@coursecodetext{}
\makeatother

% -----------------------------------------------------------------------------
% Beamer theme assignments (only apply under Beamer).
%
% We detect Beamer via \@ifclassloaded{beamer}, which is the documented,
% non-internal way. This must be wrapped in \makeatletter/\makeatother
% because \@ifclassloaded contains an @-catcode character.
% -----------------------------------------------------------------------------
\makeatletter
\@ifclassloaded{beamer}{%
  \usefonttheme{professionalfonts}
  \setbeamercolor{structure}{fg=primary-blue}
  \setbeamercolor{frametitle}{fg=primary-blue}
  \setbeamercolor{title}{fg=primary-blue}
  \setbeamercolor{subtitle}{fg=primary-gold}
  \setbeamercolor{author}{fg=primary-blue}
  \setbeamercolor{institute}{fg=neutral}
  \setbeamercolor{date}{fg=primary-gold}
  \setbeamercolor{itemize item}{fg=primary-blue}
  \setbeamercolor{itemize subitem}{fg=primary-gold}
  \setbeamercolor{alerted text}{fg=primary-gold}
  \setbeamercolor{block title}{fg=white, bg=primary-blue}
  \setbeamercolor{block body}{bg=light-bg}
  % Semantic block variants: block=definition/concept (blue), exampleblock=
  % worked example (green/positive), alertblock=key takeaway or caution (gold).
  % Keep to <=2 colored boxes per slide across ALL three variants (INV-7).
  \setbeamercolor{block title example}{fg=white, bg=positive}
  \setbeamercolor{block body example}{bg=light-bg}
  \setbeamercolor{block title alerted}{fg=white, bg=primary-gold}
  \setbeamercolor{block body alerted}{bg=light-bg}

  % ---------------------------------------------------------------------------
  % Formal title page: serif (Libertine, via \rmfamily) for title-page
  % elements only. Body text, itemize, and frametitle are intentionally left
  % on Beamer's sans-serif default -- \usebeamerfont{title} is also what
  % \transitionslide (below) uses for its big centered text, so section
  % transitions pick up the same serif treatment for free.
  % ---------------------------------------------------------------------------
  \setbeamerfont{title}{family=\rmfamily, series=\bfseries, size=\Large}
  \setbeamerfont{subtitle}{family=\rmfamily, shape=\itshape, size=\normalsize}
  \setbeamerfont{author}{family=\rmfamily, size=\normalsize}
  \setbeamerfont{institute}{family=\rmfamily, size=\scriptsize}
  \setbeamerfont{date}{family=\rmfamily, size=\scriptsize}

  \setbeamertemplate{title page}{
    \vfill
    \begin{center}
      {\usebeamerfont{title}\usebeamercolor[fg]{title}\inserttitle\par}
      \vspace{0.15cm}
      {\usebeamerfont{subtitle}\usebeamercolor[fg]{subtitle}\insertsubtitle\par}
      \vspace{0.45cm}
      {\color{primary-gold}\rule{0.42\textwidth}{0.6pt}\par}
      \vspace{0.45cm}
      {\usebeamerfont{author}\usebeamercolor[fg]{author}\insertauthor\par}
      \vspace{0.35cm}
      {\usebeamerfont{institute}\usebeamercolor[fg]{institute}\insertinstitute\par}
      \vspace{0.35cm}
      {\usebeamerfont{date}\usebeamercolor[fg]{date}\insertdate\par}
    \end{center}
    \vfill
  }

  % Minimal footer: course code (left) + page X / N (right), no navigation clutter
  \setbeamertemplate{navigation symbols}{}
  \setbeamertemplate{footline}{%
    \color{neutral}\footnotesize\hspace{0.7em}\@coursecodetext\hfill
    \insertframenumber/\inserttotalframenumber\hspace{0.7em}\vspace{0.3em}}
}{}
\makeatother

% -----------------------------------------------------------------------------
% TikZ setup — libraries the snippet gallery uses
% -----------------------------------------------------------------------------
\usepackage{tikz}
\usetikzlibrary{arrows.meta, positioning, calc, decorations.pathreplacing,
                fit, shapes.geometric, backgrounds}

% Shared TikZ styles — snippets and hand-written diagrams can pick these up
% via `\begin{tikzpicture}[every node/.style={...}]` or by naming specific
% styles (e.g., `\node[dag-node] (X) at (0,0) {X};`).
\tikzset{
  % Node styles for causal DAGs and similar
  dag-node/.style = {
    draw=primary-blue, fill=light-bg, rounded corners=2pt,
    minimum width=1.2cm, minimum height=0.9cm, align=center, font=\small
  },
  decision-node/.style = {
    draw=primary-gold, fill=white, diamond, aspect=1.8,
    minimum width=1.8cm, minimum height=1.2cm, align=center, font=\small,
    inner sep=1pt
  },
  % Edge styles — observed vs counterfactual
  observed-edge/.style = {-{Latex[length=2.5mm]}, thick, draw=jet},
  counterfactual-edge/.style = {-{Latex[length=2.5mm]}, thick, draw=neutral, dashed},
  confound-edge/.style = {-{Latex[length=2.5mm]}, thick, draw=negative, dashed},
  % Data-point styles for plots
  observed-dot/.style = {fill=primary-blue, draw=primary-blue, circle, inner sep=1.2pt},
  counterfactual-dot/.style = {fill=white, draw=neutral, circle, inner sep=1.2pt}
}

% -----------------------------------------------------------------------------
% Convenience macros
% -----------------------------------------------------------------------------
\newcommand{\muted}[1]{\textcolor{neutral}{#1}}
\newcommand{\key}[1]{\textcolor{primary-gold}{\textbf{#1}}}
\newcommand{\good}[1]{\textcolor{positive}{#1}}
\newcommand{\bad}[1]{\textcolor{negative}{#1}}

% Standout frame macro: full-bleed dark background for section transitions.
% Beamer-only (uses \begin{frame}/\setbeamercolor/\usebeamerfont) -- do not
% call from an article-class file (e.g. Notes/<CODE>/*-notes.tex). \muted,
% \key, \good, \bad, and \coursecode above are plain xcolor/gdef and are
% safe to use from either class; verified when Notes/article-class support
% was added.
% Expand: \transitionslide{Your transition title}
\newcommand{\transitionslide}[1]{%
  \begingroup
  \setbeamercolor{background canvas}{bg=primary-blue}
  \begin{frame}[plain]
    \centering
    \vfill
    {\usebeamerfont{title}\color{white}\Large #1\par}
    \vfill
  \end{frame}
  \endgroup
}
, before \begin{document})
% via \coursecode{CS401}. Shown in the footer on every slide so a deck's course
% is identifiable without opening the title page. Empty by default.
% -----------------------------------------------------------------------------
\makeatletter
\newcommand{\coursecode}[1]{\gdef\@coursecodetext{#1}}
\gdef\@coursecodetext{}
\makeatother

% -----------------------------------------------------------------------------
% Beamer theme assignments (only apply under Beamer).
%
% We detect Beamer via \@ifclassloaded{beamer}, which is the documented,
% non-internal way. This must be wrapped in \makeatletter/\makeatother
% because \@ifclassloaded contains an @-catcode character.
% -----------------------------------------------------------------------------
\makeatletter
\@ifclassloaded{beamer}{%
  \usefonttheme{professionalfonts}
  \setbeamercolor{structure}{fg=primary-blue}
  \setbeamercolor{frametitle}{fg=primary-blue}
  \setbeamercolor{title}{fg=primary-blue}
  \setbeamercolor{subtitle}{fg=primary-gold}
  \setbeamercolor{author}{fg=primary-blue}
  \setbeamercolor{institute}{fg=neutral}
  \setbeamercolor{date}{fg=primary-gold}
  \setbeamercolor{itemize item}{fg=primary-blue}
  \setbeamercolor{itemize subitem}{fg=primary-gold}
  \setbeamercolor{alerted text}{fg=primary-gold}
  \setbeamercolor{block title}{fg=white, bg=primary-blue}
  \setbeamercolor{block body}{bg=light-bg}
  % Semantic block variants: block=definition/concept (blue), exampleblock=
  % worked example (green/positive), alertblock=key takeaway or caution (gold).
  % Keep to <=2 colored boxes per slide across ALL three variants (INV-7).
  \setbeamercolor{block title example}{fg=white, bg=positive}
  \setbeamercolor{block body example}{bg=light-bg}
  \setbeamercolor{block title alerted}{fg=white, bg=primary-gold}
  \setbeamercolor{block body alerted}{bg=light-bg}

  % ---------------------------------------------------------------------------
  % Formal title page: serif (Libertine, via \rmfamily) for title-page
  % elements only. Body text, itemize, and frametitle are intentionally left
  % on Beamer's sans-serif default -- \usebeamerfont{title} is also what
  % \transitionslide (below) uses for its big centered text, so section
  % transitions pick up the same serif treatment for free.
  % ---------------------------------------------------------------------------
  \setbeamerfont{title}{family=\rmfamily, series=\bfseries, size=\Large}
  \setbeamerfont{subtitle}{family=\rmfamily, shape=\itshape, size=\normalsize}
  \setbeamerfont{author}{family=\rmfamily, size=\normalsize}
  \setbeamerfont{institute}{family=\rmfamily, size=\scriptsize}
  \setbeamerfont{date}{family=\rmfamily, size=\scriptsize}

  \setbeamertemplate{title page}{
    \vfill
    \begin{center}
      {\usebeamerfont{title}\usebeamercolor[fg]{title}\inserttitle\par}
      \vspace{0.15cm}
      {\usebeamerfont{subtitle}\usebeamercolor[fg]{subtitle}\insertsubtitle\par}
      \vspace{0.45cm}
      {\color{primary-gold}\rule{0.42\textwidth}{0.6pt}\par}
      \vspace{0.45cm}
      {\usebeamerfont{author}\usebeamercolor[fg]{author}\insertauthor\par}
      \vspace{0.35cm}
      {\usebeamerfont{institute}\usebeamercolor[fg]{institute}\insertinstitute\par}
      \vspace{0.35cm}
      {\usebeamerfont{date}\usebeamercolor[fg]{date}\insertdate\par}
    \end{center}
    \vfill
  }

  % Minimal footer: course code (left) + page X / N (right), no navigation clutter
  \setbeamertemplate{navigation symbols}{}
  \setbeamertemplate{footline}{%
    \color{neutral}\footnotesize\hspace{0.7em}\@coursecodetext\hfill
    \insertframenumber/\inserttotalframenumber\hspace{0.7em}\vspace{0.3em}}
}{}
\makeatother

% -----------------------------------------------------------------------------
% TikZ setup — libraries the snippet gallery uses
% -----------------------------------------------------------------------------
\usepackage{tikz}
\usetikzlibrary{arrows.meta, positioning, calc, decorations.pathreplacing,
                fit, shapes.geometric, backgrounds}

% Shared TikZ styles — snippets and hand-written diagrams can pick these up
% via `\begin{tikzpicture}[every node/.style={...}]` or by naming specific
% styles (e.g., `\node[dag-node] (X) at (0,0) {X};`).
\tikzset{
  % Node styles for causal DAGs and similar
  dag-node/.style = {
    draw=primary-blue, fill=light-bg, rounded corners=2pt,
    minimum width=1.2cm, minimum height=0.9cm, align=center, font=\small
  },
  decision-node/.style = {
    draw=primary-gold, fill=white, diamond, aspect=1.8,
    minimum width=1.8cm, minimum height=1.2cm, align=center, font=\small,
    inner sep=1pt
  },
  % Edge styles — observed vs counterfactual
  observed-edge/.style = {-{Latex[length=2.5mm]}, thick, draw=jet},
  counterfactual-edge/.style = {-{Latex[length=2.5mm]}, thick, draw=neutral, dashed},
  confound-edge/.style = {-{Latex[length=2.5mm]}, thick, draw=negative, dashed},
  % Data-point styles for plots
  observed-dot/.style = {fill=primary-blue, draw=primary-blue, circle, inner sep=1.2pt},
  counterfactual-dot/.style = {fill=white, draw=neutral, circle, inner sep=1.2pt}
}

% -----------------------------------------------------------------------------
% Convenience macros
% -----------------------------------------------------------------------------
\newcommand{\muted}[1]{\textcolor{neutral}{#1}}
\newcommand{\key}[1]{\textcolor{primary-gold}{\textbf{#1}}}
\newcommand{\good}[1]{\textcolor{positive}{#1}}
\newcommand{\bad}[1]{\textcolor{negative}{#1}}

% Standout frame macro: full-bleed dark background for section transitions.
% Beamer-only (uses \begin{frame}/\setbeamercolor/\usebeamerfont) -- do not
% call from an article-class file (e.g. Notes/<CODE>/*-notes.tex). \muted,
% \key, \good, \bad, and \coursecode above are plain xcolor/gdef and are
% safe to use from either class; verified when Notes/article-class support
% was added.
% Expand: \transitionslide{Your transition title}
\newcommand{\transitionslide}[1]{%
  \begingroup
  \setbeamercolor{background canvas}{bg=primary-blue}
  \begin{frame}[plain]
    \centering
    \vfill
    {\usebeamerfont{title}\color{white}\Large #1\par}
    \vfill
  \end{frame}
  \endgroup
}

% and compiling with TEXINPUTS=../Preambles:$TEXINPUTS (the /compile-latex
% skill does this automatically).
%
% PALETTE CONTRACT. Color names defined here MUST match the names in
% Quarto/theme-template.scss so Beamer and Quarto renderings use the same
% palette. The scripts/check-palette-sync.sh script enforces this.
%
% To customize: edit the HEX values below AND the matching $variable in
% Quarto/theme-template.scss, then run scripts/check-palette-sync.sh.
%
% TITLE PAGE. A formal, serif (Libertine) title-page template is defined in
% the Beamer-only block below (\setbeamertemplate{title page}) -- title,
% subtitle, a thin gold rule, then \author/\institute/\date. Frametitle and
% body text remain on Beamer's sans-serif default. No SCSS counterpart is
% needed for this (font/layout only, no new colors).
% =============================================================================

% -----------------------------------------------------------------------------
% Engine detection + encoding
%
% Our /compile-latex skill uses XeLaTeX, where inputenc is unnecessary and
% can produce encoding warnings. Only load inputenc under pdfTeX.
% -----------------------------------------------------------------------------
\usepackage{iftex}
\ifPDFTeX
  \usepackage[utf8]{inputenc}
\fi

\usepackage{amsmath, amssymb, amsthm}
\usepackage{booktabs}
\usepackage{graphicx}

% -----------------------------------------------------------------------------
% Serif face for the title page (formal/academic look). Linux Libertine is
% XeLaTeX- and pdfTeX-safe and redefines \rmdefault, so \rmfamily picks it up
% everywhere \rmfamily is used explicitly. Body text and frametitles stay on
% Beamer's own sans-serif default (unaffected) -- only the title-page
% \setbeamerfont{...} overrides below opt in to \rmfamily.
% -----------------------------------------------------------------------------
\usepackage{libertine}

% -----------------------------------------------------------------------------
% PALETTE — mirrors Quarto/theme-template.scss variable names.
% Load xcolor and define colors BEFORE anything that references them
% (notably hyperref, hypersetup, and the Beamer color assignments).
% -----------------------------------------------------------------------------
\usepackage{xcolor}

\definecolor{primary-blue}{HTML}{012169}      % headings, accents
\definecolor{primary-gold}{HTML}{B9975B}      % emphasis, borders
\definecolor{highlight-yellow}{HTML}{F2A900}  % markers, alerts
\definecolor{light-bg}{HTML}{E8EDF5}          % title / section backgrounds
\definecolor{jet}{HTML}{1A1A1A}               % body text

% Semantic colors (used in diagrams and annotations)
\definecolor{positive}{HTML}{15803D}          % good / observed / identified
\definecolor{negative}{HTML}{B91C1C}          % bad / problematic / bias
\definecolor{neutral}{HTML}{525252}           % reference / context

% Highlight accents (match .hi-* classes in the SCSS)
\definecolor{hi-slate}{HTML}{314F4F}
\definecolor{hi-green}{HTML}{2E8B57}
\definecolor{hi-red}{HTML}{C41E3A}

% Semantic aliases so skill templates and snippets can write
% \textcolor{accent}{...} without hard-coding "primary-blue".
\colorlet{accent}{primary-blue}
\colorlet{accent2}{primary-gold}

% -----------------------------------------------------------------------------
% Hyperref — loaded AFTER xcolor + palette so link colors resolve.
% -----------------------------------------------------------------------------
\usepackage{hyperref}
\hypersetup{colorlinks=true, linkcolor=primary-blue, urlcolor=primary-blue,
            citecolor=primary-blue, pdfborder={0 0 0}}

% -----------------------------------------------------------------------------
% Course code — set once per deck (after % =============================================================================
% Preambles/header.tex — shared LaTeX/Beamer preamble for this project
%
% Use from a Beamer lecture by prepending:
%     % =============================================================================
% Preambles/header.tex — shared LaTeX/Beamer preamble for this project
%
% Use from a Beamer lecture by prepending:
%     \input{header}
% and compiling with TEXINPUTS=../Preambles:$TEXINPUTS (the /compile-latex
% skill does this automatically).
%
% PALETTE CONTRACT. Color names defined here MUST match the names in
% Quarto/theme-template.scss so Beamer and Quarto renderings use the same
% palette. The scripts/check-palette-sync.sh script enforces this.
%
% To customize: edit the HEX values below AND the matching $variable in
% Quarto/theme-template.scss, then run scripts/check-palette-sync.sh.
%
% TITLE PAGE. A formal, serif (Libertine) title-page template is defined in
% the Beamer-only block below (\setbeamertemplate{title page}) -- title,
% subtitle, a thin gold rule, then \author/\institute/\date. Frametitle and
% body text remain on Beamer's sans-serif default. No SCSS counterpart is
% needed for this (font/layout only, no new colors).
% =============================================================================

% -----------------------------------------------------------------------------
% Engine detection + encoding
%
% Our /compile-latex skill uses XeLaTeX, where inputenc is unnecessary and
% can produce encoding warnings. Only load inputenc under pdfTeX.
% -----------------------------------------------------------------------------
\usepackage{iftex}
\ifPDFTeX
  \usepackage[utf8]{inputenc}
\fi

\usepackage{amsmath, amssymb, amsthm}
\usepackage{booktabs}
\usepackage{graphicx}

% -----------------------------------------------------------------------------
% Serif face for the title page (formal/academic look). Linux Libertine is
% XeLaTeX- and pdfTeX-safe and redefines \rmdefault, so \rmfamily picks it up
% everywhere \rmfamily is used explicitly. Body text and frametitles stay on
% Beamer's own sans-serif default (unaffected) -- only the title-page
% \setbeamerfont{...} overrides below opt in to \rmfamily.
% -----------------------------------------------------------------------------
\usepackage{libertine}

% -----------------------------------------------------------------------------
% PALETTE — mirrors Quarto/theme-template.scss variable names.
% Load xcolor and define colors BEFORE anything that references them
% (notably hyperref, hypersetup, and the Beamer color assignments).
% -----------------------------------------------------------------------------
\usepackage{xcolor}

\definecolor{primary-blue}{HTML}{012169}      % headings, accents
\definecolor{primary-gold}{HTML}{B9975B}      % emphasis, borders
\definecolor{highlight-yellow}{HTML}{F2A900}  % markers, alerts
\definecolor{light-bg}{HTML}{E8EDF5}          % title / section backgrounds
\definecolor{jet}{HTML}{1A1A1A}               % body text

% Semantic colors (used in diagrams and annotations)
\definecolor{positive}{HTML}{15803D}          % good / observed / identified
\definecolor{negative}{HTML}{B91C1C}          % bad / problematic / bias
\definecolor{neutral}{HTML}{525252}           % reference / context

% Highlight accents (match .hi-* classes in the SCSS)
\definecolor{hi-slate}{HTML}{314F4F}
\definecolor{hi-green}{HTML}{2E8B57}
\definecolor{hi-red}{HTML}{C41E3A}

% Semantic aliases so skill templates and snippets can write
% \textcolor{accent}{...} without hard-coding "primary-blue".
\colorlet{accent}{primary-blue}
\colorlet{accent2}{primary-gold}

% -----------------------------------------------------------------------------
% Hyperref — loaded AFTER xcolor + palette so link colors resolve.
% -----------------------------------------------------------------------------
\usepackage{hyperref}
\hypersetup{colorlinks=true, linkcolor=primary-blue, urlcolor=primary-blue,
            citecolor=primary-blue, pdfborder={0 0 0}}

% -----------------------------------------------------------------------------
% Course code — set once per deck (after \input{header}, before \begin{document})
% via \coursecode{CS401}. Shown in the footer on every slide so a deck's course
% is identifiable without opening the title page. Empty by default.
% -----------------------------------------------------------------------------
\makeatletter
\newcommand{\coursecode}[1]{\gdef\@coursecodetext{#1}}
\gdef\@coursecodetext{}
\makeatother

% -----------------------------------------------------------------------------
% Beamer theme assignments (only apply under Beamer).
%
% We detect Beamer via \@ifclassloaded{beamer}, which is the documented,
% non-internal way. This must be wrapped in \makeatletter/\makeatother
% because \@ifclassloaded contains an @-catcode character.
% -----------------------------------------------------------------------------
\makeatletter
\@ifclassloaded{beamer}{%
  \usefonttheme{professionalfonts}
  \setbeamercolor{structure}{fg=primary-blue}
  \setbeamercolor{frametitle}{fg=primary-blue}
  \setbeamercolor{title}{fg=primary-blue}
  \setbeamercolor{subtitle}{fg=primary-gold}
  \setbeamercolor{author}{fg=primary-blue}
  \setbeamercolor{institute}{fg=neutral}
  \setbeamercolor{date}{fg=primary-gold}
  \setbeamercolor{itemize item}{fg=primary-blue}
  \setbeamercolor{itemize subitem}{fg=primary-gold}
  \setbeamercolor{alerted text}{fg=primary-gold}
  \setbeamercolor{block title}{fg=white, bg=primary-blue}
  \setbeamercolor{block body}{bg=light-bg}
  % Semantic block variants: block=definition/concept (blue), exampleblock=
  % worked example (green/positive), alertblock=key takeaway or caution (gold).
  % Keep to <=2 colored boxes per slide across ALL three variants (INV-7).
  \setbeamercolor{block title example}{fg=white, bg=positive}
  \setbeamercolor{block body example}{bg=light-bg}
  \setbeamercolor{block title alerted}{fg=white, bg=primary-gold}
  \setbeamercolor{block body alerted}{bg=light-bg}

  % ---------------------------------------------------------------------------
  % Formal title page: serif (Libertine, via \rmfamily) for title-page
  % elements only. Body text, itemize, and frametitle are intentionally left
  % on Beamer's sans-serif default -- \usebeamerfont{title} is also what
  % \transitionslide (below) uses for its big centered text, so section
  % transitions pick up the same serif treatment for free.
  % ---------------------------------------------------------------------------
  \setbeamerfont{title}{family=\rmfamily, series=\bfseries, size=\Large}
  \setbeamerfont{subtitle}{family=\rmfamily, shape=\itshape, size=\normalsize}
  \setbeamerfont{author}{family=\rmfamily, size=\normalsize}
  \setbeamerfont{institute}{family=\rmfamily, size=\scriptsize}
  \setbeamerfont{date}{family=\rmfamily, size=\scriptsize}

  \setbeamertemplate{title page}{
    \vfill
    \begin{center}
      {\usebeamerfont{title}\usebeamercolor[fg]{title}\inserttitle\par}
      \vspace{0.15cm}
      {\usebeamerfont{subtitle}\usebeamercolor[fg]{subtitle}\insertsubtitle\par}
      \vspace{0.45cm}
      {\color{primary-gold}\rule{0.42\textwidth}{0.6pt}\par}
      \vspace{0.45cm}
      {\usebeamerfont{author}\usebeamercolor[fg]{author}\insertauthor\par}
      \vspace{0.35cm}
      {\usebeamerfont{institute}\usebeamercolor[fg]{institute}\insertinstitute\par}
      \vspace{0.35cm}
      {\usebeamerfont{date}\usebeamercolor[fg]{date}\insertdate\par}
    \end{center}
    \vfill
  }

  % Minimal footer: course code (left) + page X / N (right), no navigation clutter
  \setbeamertemplate{navigation symbols}{}
  \setbeamertemplate{footline}{%
    \color{neutral}\footnotesize\hspace{0.7em}\@coursecodetext\hfill
    \insertframenumber/\inserttotalframenumber\hspace{0.7em}\vspace{0.3em}}
}{}
\makeatother

% -----------------------------------------------------------------------------
% TikZ setup — libraries the snippet gallery uses
% -----------------------------------------------------------------------------
\usepackage{tikz}
\usetikzlibrary{arrows.meta, positioning, calc, decorations.pathreplacing,
                fit, shapes.geometric, backgrounds}

% Shared TikZ styles — snippets and hand-written diagrams can pick these up
% via `\begin{tikzpicture}[every node/.style={...}]` or by naming specific
% styles (e.g., `\node[dag-node] (X) at (0,0) {X};`).
\tikzset{
  % Node styles for causal DAGs and similar
  dag-node/.style = {
    draw=primary-blue, fill=light-bg, rounded corners=2pt,
    minimum width=1.2cm, minimum height=0.9cm, align=center, font=\small
  },
  decision-node/.style = {
    draw=primary-gold, fill=white, diamond, aspect=1.8,
    minimum width=1.8cm, minimum height=1.2cm, align=center, font=\small,
    inner sep=1pt
  },
  % Edge styles — observed vs counterfactual
  observed-edge/.style = {-{Latex[length=2.5mm]}, thick, draw=jet},
  counterfactual-edge/.style = {-{Latex[length=2.5mm]}, thick, draw=neutral, dashed},
  confound-edge/.style = {-{Latex[length=2.5mm]}, thick, draw=negative, dashed},
  % Data-point styles for plots
  observed-dot/.style = {fill=primary-blue, draw=primary-blue, circle, inner sep=1.2pt},
  counterfactual-dot/.style = {fill=white, draw=neutral, circle, inner sep=1.2pt}
}

% -----------------------------------------------------------------------------
% Convenience macros
% -----------------------------------------------------------------------------
\newcommand{\muted}[1]{\textcolor{neutral}{#1}}
\newcommand{\key}[1]{\textcolor{primary-gold}{\textbf{#1}}}
\newcommand{\good}[1]{\textcolor{positive}{#1}}
\newcommand{\bad}[1]{\textcolor{negative}{#1}}

% Standout frame macro: full-bleed dark background for section transitions.
% Beamer-only (uses \begin{frame}/\setbeamercolor/\usebeamerfont) -- do not
% call from an article-class file (e.g. Notes/<CODE>/*-notes.tex). \muted,
% \key, \good, \bad, and \coursecode above are plain xcolor/gdef and are
% safe to use from either class; verified when Notes/article-class support
% was added.
% Expand: \transitionslide{Your transition title}
\newcommand{\transitionslide}[1]{%
  \begingroup
  \setbeamercolor{background canvas}{bg=primary-blue}
  \begin{frame}[plain]
    \centering
    \vfill
    {\usebeamerfont{title}\color{white}\Large #1\par}
    \vfill
  \end{frame}
  \endgroup
}

% and compiling with TEXINPUTS=../Preambles:$TEXINPUTS (the /compile-latex
% skill does this automatically).
%
% PALETTE CONTRACT. Color names defined here MUST match the names in
% Quarto/theme-template.scss so Beamer and Quarto renderings use the same
% palette. The scripts/check-palette-sync.sh script enforces this.
%
% To customize: edit the HEX values below AND the matching $variable in
% Quarto/theme-template.scss, then run scripts/check-palette-sync.sh.
%
% TITLE PAGE. A formal, serif (Libertine) title-page template is defined in
% the Beamer-only block below (\setbeamertemplate{title page}) -- title,
% subtitle, a thin gold rule, then \author/\institute/\date. Frametitle and
% body text remain on Beamer's sans-serif default. No SCSS counterpart is
% needed for this (font/layout only, no new colors).
% =============================================================================

% -----------------------------------------------------------------------------
% Engine detection + encoding
%
% Our /compile-latex skill uses XeLaTeX, where inputenc is unnecessary and
% can produce encoding warnings. Only load inputenc under pdfTeX.
% -----------------------------------------------------------------------------
\usepackage{iftex}
\ifPDFTeX
  \usepackage[utf8]{inputenc}
\fi

\usepackage{amsmath, amssymb, amsthm}
\usepackage{booktabs}
\usepackage{graphicx}

% -----------------------------------------------------------------------------
% Serif face for the title page (formal/academic look). Linux Libertine is
% XeLaTeX- and pdfTeX-safe and redefines \rmdefault, so \rmfamily picks it up
% everywhere \rmfamily is used explicitly. Body text and frametitles stay on
% Beamer's own sans-serif default (unaffected) -- only the title-page
% \setbeamerfont{...} overrides below opt in to \rmfamily.
% -----------------------------------------------------------------------------
\usepackage{libertine}

% -----------------------------------------------------------------------------
% PALETTE — mirrors Quarto/theme-template.scss variable names.
% Load xcolor and define colors BEFORE anything that references them
% (notably hyperref, hypersetup, and the Beamer color assignments).
% -----------------------------------------------------------------------------
\usepackage{xcolor}

\definecolor{primary-blue}{HTML}{012169}      % headings, accents
\definecolor{primary-gold}{HTML}{B9975B}      % emphasis, borders
\definecolor{highlight-yellow}{HTML}{F2A900}  % markers, alerts
\definecolor{light-bg}{HTML}{E8EDF5}          % title / section backgrounds
\definecolor{jet}{HTML}{1A1A1A}               % body text

% Semantic colors (used in diagrams and annotations)
\definecolor{positive}{HTML}{15803D}          % good / observed / identified
\definecolor{negative}{HTML}{B91C1C}          % bad / problematic / bias
\definecolor{neutral}{HTML}{525252}           % reference / context

% Highlight accents (match .hi-* classes in the SCSS)
\definecolor{hi-slate}{HTML}{314F4F}
\definecolor{hi-green}{HTML}{2E8B57}
\definecolor{hi-red}{HTML}{C41E3A}

% Semantic aliases so skill templates and snippets can write
% \textcolor{accent}{...} without hard-coding "primary-blue".
\colorlet{accent}{primary-blue}
\colorlet{accent2}{primary-gold}

% -----------------------------------------------------------------------------
% Hyperref — loaded AFTER xcolor + palette so link colors resolve.
% -----------------------------------------------------------------------------
\usepackage{hyperref}
\hypersetup{colorlinks=true, linkcolor=primary-blue, urlcolor=primary-blue,
            citecolor=primary-blue, pdfborder={0 0 0}}

% -----------------------------------------------------------------------------
% Course code — set once per deck (after % =============================================================================
% Preambles/header.tex — shared LaTeX/Beamer preamble for this project
%
% Use from a Beamer lecture by prepending:
%     \input{header}
% and compiling with TEXINPUTS=../Preambles:$TEXINPUTS (the /compile-latex
% skill does this automatically).
%
% PALETTE CONTRACT. Color names defined here MUST match the names in
% Quarto/theme-template.scss so Beamer and Quarto renderings use the same
% palette. The scripts/check-palette-sync.sh script enforces this.
%
% To customize: edit the HEX values below AND the matching $variable in
% Quarto/theme-template.scss, then run scripts/check-palette-sync.sh.
%
% TITLE PAGE. A formal, serif (Libertine) title-page template is defined in
% the Beamer-only block below (\setbeamertemplate{title page}) -- title,
% subtitle, a thin gold rule, then \author/\institute/\date. Frametitle and
% body text remain on Beamer's sans-serif default. No SCSS counterpart is
% needed for this (font/layout only, no new colors).
% =============================================================================

% -----------------------------------------------------------------------------
% Engine detection + encoding
%
% Our /compile-latex skill uses XeLaTeX, where inputenc is unnecessary and
% can produce encoding warnings. Only load inputenc under pdfTeX.
% -----------------------------------------------------------------------------
\usepackage{iftex}
\ifPDFTeX
  \usepackage[utf8]{inputenc}
\fi

\usepackage{amsmath, amssymb, amsthm}
\usepackage{booktabs}
\usepackage{graphicx}

% -----------------------------------------------------------------------------
% Serif face for the title page (formal/academic look). Linux Libertine is
% XeLaTeX- and pdfTeX-safe and redefines \rmdefault, so \rmfamily picks it up
% everywhere \rmfamily is used explicitly. Body text and frametitles stay on
% Beamer's own sans-serif default (unaffected) -- only the title-page
% \setbeamerfont{...} overrides below opt in to \rmfamily.
% -----------------------------------------------------------------------------
\usepackage{libertine}

% -----------------------------------------------------------------------------
% PALETTE — mirrors Quarto/theme-template.scss variable names.
% Load xcolor and define colors BEFORE anything that references them
% (notably hyperref, hypersetup, and the Beamer color assignments).
% -----------------------------------------------------------------------------
\usepackage{xcolor}

\definecolor{primary-blue}{HTML}{012169}      % headings, accents
\definecolor{primary-gold}{HTML}{B9975B}      % emphasis, borders
\definecolor{highlight-yellow}{HTML}{F2A900}  % markers, alerts
\definecolor{light-bg}{HTML}{E8EDF5}          % title / section backgrounds
\definecolor{jet}{HTML}{1A1A1A}               % body text

% Semantic colors (used in diagrams and annotations)
\definecolor{positive}{HTML}{15803D}          % good / observed / identified
\definecolor{negative}{HTML}{B91C1C}          % bad / problematic / bias
\definecolor{neutral}{HTML}{525252}           % reference / context

% Highlight accents (match .hi-* classes in the SCSS)
\definecolor{hi-slate}{HTML}{314F4F}
\definecolor{hi-green}{HTML}{2E8B57}
\definecolor{hi-red}{HTML}{C41E3A}

% Semantic aliases so skill templates and snippets can write
% \textcolor{accent}{...} without hard-coding "primary-blue".
\colorlet{accent}{primary-blue}
\colorlet{accent2}{primary-gold}

% -----------------------------------------------------------------------------
% Hyperref — loaded AFTER xcolor + palette so link colors resolve.
% -----------------------------------------------------------------------------
\usepackage{hyperref}
\hypersetup{colorlinks=true, linkcolor=primary-blue, urlcolor=primary-blue,
            citecolor=primary-blue, pdfborder={0 0 0}}

% -----------------------------------------------------------------------------
% Course code — set once per deck (after \input{header}, before \begin{document})
% via \coursecode{CS401}. Shown in the footer on every slide so a deck's course
% is identifiable without opening the title page. Empty by default.
% -----------------------------------------------------------------------------
\makeatletter
\newcommand{\coursecode}[1]{\gdef\@coursecodetext{#1}}
\gdef\@coursecodetext{}
\makeatother

% -----------------------------------------------------------------------------
% Beamer theme assignments (only apply under Beamer).
%
% We detect Beamer via \@ifclassloaded{beamer}, which is the documented,
% non-internal way. This must be wrapped in \makeatletter/\makeatother
% because \@ifclassloaded contains an @-catcode character.
% -----------------------------------------------------------------------------
\makeatletter
\@ifclassloaded{beamer}{%
  \usefonttheme{professionalfonts}
  \setbeamercolor{structure}{fg=primary-blue}
  \setbeamercolor{frametitle}{fg=primary-blue}
  \setbeamercolor{title}{fg=primary-blue}
  \setbeamercolor{subtitle}{fg=primary-gold}
  \setbeamercolor{author}{fg=primary-blue}
  \setbeamercolor{institute}{fg=neutral}
  \setbeamercolor{date}{fg=primary-gold}
  \setbeamercolor{itemize item}{fg=primary-blue}
  \setbeamercolor{itemize subitem}{fg=primary-gold}
  \setbeamercolor{alerted text}{fg=primary-gold}
  \setbeamercolor{block title}{fg=white, bg=primary-blue}
  \setbeamercolor{block body}{bg=light-bg}
  % Semantic block variants: block=definition/concept (blue), exampleblock=
  % worked example (green/positive), alertblock=key takeaway or caution (gold).
  % Keep to <=2 colored boxes per slide across ALL three variants (INV-7).
  \setbeamercolor{block title example}{fg=white, bg=positive}
  \setbeamercolor{block body example}{bg=light-bg}
  \setbeamercolor{block title alerted}{fg=white, bg=primary-gold}
  \setbeamercolor{block body alerted}{bg=light-bg}

  % ---------------------------------------------------------------------------
  % Formal title page: serif (Libertine, via \rmfamily) for title-page
  % elements only. Body text, itemize, and frametitle are intentionally left
  % on Beamer's sans-serif default -- \usebeamerfont{title} is also what
  % \transitionslide (below) uses for its big centered text, so section
  % transitions pick up the same serif treatment for free.
  % ---------------------------------------------------------------------------
  \setbeamerfont{title}{family=\rmfamily, series=\bfseries, size=\Large}
  \setbeamerfont{subtitle}{family=\rmfamily, shape=\itshape, size=\normalsize}
  \setbeamerfont{author}{family=\rmfamily, size=\normalsize}
  \setbeamerfont{institute}{family=\rmfamily, size=\scriptsize}
  \setbeamerfont{date}{family=\rmfamily, size=\scriptsize}

  \setbeamertemplate{title page}{
    \vfill
    \begin{center}
      {\usebeamerfont{title}\usebeamercolor[fg]{title}\inserttitle\par}
      \vspace{0.15cm}
      {\usebeamerfont{subtitle}\usebeamercolor[fg]{subtitle}\insertsubtitle\par}
      \vspace{0.45cm}
      {\color{primary-gold}\rule{0.42\textwidth}{0.6pt}\par}
      \vspace{0.45cm}
      {\usebeamerfont{author}\usebeamercolor[fg]{author}\insertauthor\par}
      \vspace{0.35cm}
      {\usebeamerfont{institute}\usebeamercolor[fg]{institute}\insertinstitute\par}
      \vspace{0.35cm}
      {\usebeamerfont{date}\usebeamercolor[fg]{date}\insertdate\par}
    \end{center}
    \vfill
  }

  % Minimal footer: course code (left) + page X / N (right), no navigation clutter
  \setbeamertemplate{navigation symbols}{}
  \setbeamertemplate{footline}{%
    \color{neutral}\footnotesize\hspace{0.7em}\@coursecodetext\hfill
    \insertframenumber/\inserttotalframenumber\hspace{0.7em}\vspace{0.3em}}
}{}
\makeatother

% -----------------------------------------------------------------------------
% TikZ setup — libraries the snippet gallery uses
% -----------------------------------------------------------------------------
\usepackage{tikz}
\usetikzlibrary{arrows.meta, positioning, calc, decorations.pathreplacing,
                fit, shapes.geometric, backgrounds}

% Shared TikZ styles — snippets and hand-written diagrams can pick these up
% via `\begin{tikzpicture}[every node/.style={...}]` or by naming specific
% styles (e.g., `\node[dag-node] (X) at (0,0) {X};`).
\tikzset{
  % Node styles for causal DAGs and similar
  dag-node/.style = {
    draw=primary-blue, fill=light-bg, rounded corners=2pt,
    minimum width=1.2cm, minimum height=0.9cm, align=center, font=\small
  },
  decision-node/.style = {
    draw=primary-gold, fill=white, diamond, aspect=1.8,
    minimum width=1.8cm, minimum height=1.2cm, align=center, font=\small,
    inner sep=1pt
  },
  % Edge styles — observed vs counterfactual
  observed-edge/.style = {-{Latex[length=2.5mm]}, thick, draw=jet},
  counterfactual-edge/.style = {-{Latex[length=2.5mm]}, thick, draw=neutral, dashed},
  confound-edge/.style = {-{Latex[length=2.5mm]}, thick, draw=negative, dashed},
  % Data-point styles for plots
  observed-dot/.style = {fill=primary-blue, draw=primary-blue, circle, inner sep=1.2pt},
  counterfactual-dot/.style = {fill=white, draw=neutral, circle, inner sep=1.2pt}
}

% -----------------------------------------------------------------------------
% Convenience macros
% -----------------------------------------------------------------------------
\newcommand{\muted}[1]{\textcolor{neutral}{#1}}
\newcommand{\key}[1]{\textcolor{primary-gold}{\textbf{#1}}}
\newcommand{\good}[1]{\textcolor{positive}{#1}}
\newcommand{\bad}[1]{\textcolor{negative}{#1}}

% Standout frame macro: full-bleed dark background for section transitions.
% Beamer-only (uses \begin{frame}/\setbeamercolor/\usebeamerfont) -- do not
% call from an article-class file (e.g. Notes/<CODE>/*-notes.tex). \muted,
% \key, \good, \bad, and \coursecode above are plain xcolor/gdef and are
% safe to use from either class; verified when Notes/article-class support
% was added.
% Expand: \transitionslide{Your transition title}
\newcommand{\transitionslide}[1]{%
  \begingroup
  \setbeamercolor{background canvas}{bg=primary-blue}
  \begin{frame}[plain]
    \centering
    \vfill
    {\usebeamerfont{title}\color{white}\Large #1\par}
    \vfill
  \end{frame}
  \endgroup
}
, before \begin{document})
% via \coursecode{CS401}. Shown in the footer on every slide so a deck's course
% is identifiable without opening the title page. Empty by default.
% -----------------------------------------------------------------------------
\makeatletter
\newcommand{\coursecode}[1]{\gdef\@coursecodetext{#1}}
\gdef\@coursecodetext{}
\makeatother

% -----------------------------------------------------------------------------
% Beamer theme assignments (only apply under Beamer).
%
% We detect Beamer via \@ifclassloaded{beamer}, which is the documented,
% non-internal way. This must be wrapped in \makeatletter/\makeatother
% because \@ifclassloaded contains an @-catcode character.
% -----------------------------------------------------------------------------
\makeatletter
\@ifclassloaded{beamer}{%
  \usefonttheme{professionalfonts}
  \setbeamercolor{structure}{fg=primary-blue}
  \setbeamercolor{frametitle}{fg=primary-blue}
  \setbeamercolor{title}{fg=primary-blue}
  \setbeamercolor{subtitle}{fg=primary-gold}
  \setbeamercolor{author}{fg=primary-blue}
  \setbeamercolor{institute}{fg=neutral}
  \setbeamercolor{date}{fg=primary-gold}
  \setbeamercolor{itemize item}{fg=primary-blue}
  \setbeamercolor{itemize subitem}{fg=primary-gold}
  \setbeamercolor{alerted text}{fg=primary-gold}
  \setbeamercolor{block title}{fg=white, bg=primary-blue}
  \setbeamercolor{block body}{bg=light-bg}
  % Semantic block variants: block=definition/concept (blue), exampleblock=
  % worked example (green/positive), alertblock=key takeaway or caution (gold).
  % Keep to <=2 colored boxes per slide across ALL three variants (INV-7).
  \setbeamercolor{block title example}{fg=white, bg=positive}
  \setbeamercolor{block body example}{bg=light-bg}
  \setbeamercolor{block title alerted}{fg=white, bg=primary-gold}
  \setbeamercolor{block body alerted}{bg=light-bg}

  % ---------------------------------------------------------------------------
  % Formal title page: serif (Libertine, via \rmfamily) for title-page
  % elements only. Body text, itemize, and frametitle are intentionally left
  % on Beamer's sans-serif default -- \usebeamerfont{title} is also what
  % \transitionslide (below) uses for its big centered text, so section
  % transitions pick up the same serif treatment for free.
  % ---------------------------------------------------------------------------
  \setbeamerfont{title}{family=\rmfamily, series=\bfseries, size=\Large}
  \setbeamerfont{subtitle}{family=\rmfamily, shape=\itshape, size=\normalsize}
  \setbeamerfont{author}{family=\rmfamily, size=\normalsize}
  \setbeamerfont{institute}{family=\rmfamily, size=\scriptsize}
  \setbeamerfont{date}{family=\rmfamily, size=\scriptsize}

  \setbeamertemplate{title page}{
    \vfill
    \begin{center}
      {\usebeamerfont{title}\usebeamercolor[fg]{title}\inserttitle\par}
      \vspace{0.15cm}
      {\usebeamerfont{subtitle}\usebeamercolor[fg]{subtitle}\insertsubtitle\par}
      \vspace{0.45cm}
      {\color{primary-gold}\rule{0.42\textwidth}{0.6pt}\par}
      \vspace{0.45cm}
      {\usebeamerfont{author}\usebeamercolor[fg]{author}\insertauthor\par}
      \vspace{0.35cm}
      {\usebeamerfont{institute}\usebeamercolor[fg]{institute}\insertinstitute\par}
      \vspace{0.35cm}
      {\usebeamerfont{date}\usebeamercolor[fg]{date}\insertdate\par}
    \end{center}
    \vfill
  }

  % Minimal footer: course code (left) + page X / N (right), no navigation clutter
  \setbeamertemplate{navigation symbols}{}
  \setbeamertemplate{footline}{%
    \color{neutral}\footnotesize\hspace{0.7em}\@coursecodetext\hfill
    \insertframenumber/\inserttotalframenumber\hspace{0.7em}\vspace{0.3em}}
}{}
\makeatother

% -----------------------------------------------------------------------------
% TikZ setup — libraries the snippet gallery uses
% -----------------------------------------------------------------------------
\usepackage{tikz}
\usetikzlibrary{arrows.meta, positioning, calc, decorations.pathreplacing,
                fit, shapes.geometric, backgrounds}

% Shared TikZ styles — snippets and hand-written diagrams can pick these up
% via `\begin{tikzpicture}[every node/.style={...}]` or by naming specific
% styles (e.g., `\node[dag-node] (X) at (0,0) {X};`).
\tikzset{
  % Node styles for causal DAGs and similar
  dag-node/.style = {
    draw=primary-blue, fill=light-bg, rounded corners=2pt,
    minimum width=1.2cm, minimum height=0.9cm, align=center, font=\small
  },
  decision-node/.style = {
    draw=primary-gold, fill=white, diamond, aspect=1.8,
    minimum width=1.8cm, minimum height=1.2cm, align=center, font=\small,
    inner sep=1pt
  },
  % Edge styles — observed vs counterfactual
  observed-edge/.style = {-{Latex[length=2.5mm]}, thick, draw=jet},
  counterfactual-edge/.style = {-{Latex[length=2.5mm]}, thick, draw=neutral, dashed},
  confound-edge/.style = {-{Latex[length=2.5mm]}, thick, draw=negative, dashed},
  % Data-point styles for plots
  observed-dot/.style = {fill=primary-blue, draw=primary-blue, circle, inner sep=1.2pt},
  counterfactual-dot/.style = {fill=white, draw=neutral, circle, inner sep=1.2pt}
}

% -----------------------------------------------------------------------------
% Convenience macros
% -----------------------------------------------------------------------------
\newcommand{\muted}[1]{\textcolor{neutral}{#1}}
\newcommand{\key}[1]{\textcolor{primary-gold}{\textbf{#1}}}
\newcommand{\good}[1]{\textcolor{positive}{#1}}
\newcommand{\bad}[1]{\textcolor{negative}{#1}}

% Standout frame macro: full-bleed dark background for section transitions.
% Beamer-only (uses \begin{frame}/\setbeamercolor/\usebeamerfont) -- do not
% call from an article-class file (e.g. Notes/<CODE>/*-notes.tex). \muted,
% \key, \good, \bad, and \coursecode above are plain xcolor/gdef and are
% safe to use from either class; verified when Notes/article-class support
% was added.
% Expand: \transitionslide{Your transition title}
\newcommand{\transitionslide}[1]{%
  \begingroup
  \setbeamercolor{background canvas}{bg=primary-blue}
  \begin{frame}[plain]
    \centering
    \vfill
    {\usebeamerfont{title}\color{white}\Large #1\par}
    \vfill
  \end{frame}
  \endgroup
}
, before \begin{document})
% via \coursecode{CS401}. Shown in the footer on every slide so a deck's course
% is identifiable without opening the title page. Empty by default.
% -----------------------------------------------------------------------------
\makeatletter
\newcommand{\coursecode}[1]{\gdef\@coursecodetext{#1}}
\gdef\@coursecodetext{}
\makeatother

% -----------------------------------------------------------------------------
% Beamer theme assignments (only apply under Beamer).
%
% We detect Beamer via \@ifclassloaded{beamer}, which is the documented,
% non-internal way. This must be wrapped in \makeatletter/\makeatother
% because \@ifclassloaded contains an @-catcode character.
% -----------------------------------------------------------------------------
\makeatletter
\@ifclassloaded{beamer}{%
  \usefonttheme{professionalfonts}
  \setbeamercolor{structure}{fg=primary-blue}
  \setbeamercolor{frametitle}{fg=primary-blue}
  \setbeamercolor{title}{fg=primary-blue}
  \setbeamercolor{subtitle}{fg=primary-gold}
  \setbeamercolor{author}{fg=primary-blue}
  \setbeamercolor{institute}{fg=neutral}
  \setbeamercolor{date}{fg=primary-gold}
  \setbeamercolor{itemize item}{fg=primary-blue}
  \setbeamercolor{itemize subitem}{fg=primary-gold}
  \setbeamercolor{alerted text}{fg=primary-gold}
  \setbeamercolor{block title}{fg=white, bg=primary-blue}
  \setbeamercolor{block body}{bg=light-bg}
  % Semantic block variants: block=definition/concept (blue), exampleblock=
  % worked example (green/positive), alertblock=key takeaway or caution (gold).
  % Keep to <=2 colored boxes per slide across ALL three variants (INV-7).
  \setbeamercolor{block title example}{fg=white, bg=positive}
  \setbeamercolor{block body example}{bg=light-bg}
  \setbeamercolor{block title alerted}{fg=white, bg=primary-gold}
  \setbeamercolor{block body alerted}{bg=light-bg}

  % ---------------------------------------------------------------------------
  % Formal title page: serif (Libertine, via \rmfamily) for title-page
  % elements only. Body text, itemize, and frametitle are intentionally left
  % on Beamer's sans-serif default -- \usebeamerfont{title} is also what
  % \transitionslide (below) uses for its big centered text, so section
  % transitions pick up the same serif treatment for free.
  % ---------------------------------------------------------------------------
  \setbeamerfont{title}{family=\rmfamily, series=\bfseries, size=\Large}
  \setbeamerfont{subtitle}{family=\rmfamily, shape=\itshape, size=\normalsize}
  \setbeamerfont{author}{family=\rmfamily, size=\normalsize}
  \setbeamerfont{institute}{family=\rmfamily, size=\scriptsize}
  \setbeamerfont{date}{family=\rmfamily, size=\scriptsize}

  \setbeamertemplate{title page}{
    \vfill
    \begin{center}
      {\usebeamerfont{title}\usebeamercolor[fg]{title}\inserttitle\par}
      \vspace{0.15cm}
      {\usebeamerfont{subtitle}\usebeamercolor[fg]{subtitle}\insertsubtitle\par}
      \vspace{0.45cm}
      {\color{primary-gold}\rule{0.42\textwidth}{0.6pt}\par}
      \vspace{0.45cm}
      {\usebeamerfont{author}\usebeamercolor[fg]{author}\insertauthor\par}
      \vspace{0.35cm}
      {\usebeamerfont{institute}\usebeamercolor[fg]{institute}\insertinstitute\par}
      \vspace{0.35cm}
      {\usebeamerfont{date}\usebeamercolor[fg]{date}\insertdate\par}
    \end{center}
    \vfill
  }

  % Minimal footer: course code (left) + page X / N (right), no navigation clutter
  \setbeamertemplate{navigation symbols}{}
  \setbeamertemplate{footline}{%
    \color{neutral}\footnotesize\hspace{0.7em}\@coursecodetext\hfill
    \insertframenumber/\inserttotalframenumber\hspace{0.7em}\vspace{0.3em}}
}{}
\makeatother

% -----------------------------------------------------------------------------
% TikZ setup — libraries the snippet gallery uses
% -----------------------------------------------------------------------------
\usepackage{tikz}
\usetikzlibrary{arrows.meta, positioning, calc, decorations.pathreplacing,
                fit, shapes.geometric, backgrounds}

% Shared TikZ styles — snippets and hand-written diagrams can pick these up
% via `\begin{tikzpicture}[every node/.style={...}]` or by naming specific
% styles (e.g., `\node[dag-node] (X) at (0,0) {X};`).
\tikzset{
  % Node styles for causal DAGs and similar
  dag-node/.style = {
    draw=primary-blue, fill=light-bg, rounded corners=2pt,
    minimum width=1.2cm, minimum height=0.9cm, align=center, font=\small
  },
  decision-node/.style = {
    draw=primary-gold, fill=white, diamond, aspect=1.8,
    minimum width=1.8cm, minimum height=1.2cm, align=center, font=\small,
    inner sep=1pt
  },
  % Edge styles — observed vs counterfactual
  observed-edge/.style = {-{Latex[length=2.5mm]}, thick, draw=jet},
  counterfactual-edge/.style = {-{Latex[length=2.5mm]}, thick, draw=neutral, dashed},
  confound-edge/.style = {-{Latex[length=2.5mm]}, thick, draw=negative, dashed},
  % Data-point styles for plots
  observed-dot/.style = {fill=primary-blue, draw=primary-blue, circle, inner sep=1.2pt},
  counterfactual-dot/.style = {fill=white, draw=neutral, circle, inner sep=1.2pt}
}

% -----------------------------------------------------------------------------
% Convenience macros
% -----------------------------------------------------------------------------
\newcommand{\muted}[1]{\textcolor{neutral}{#1}}
\newcommand{\key}[1]{\textcolor{primary-gold}{\textbf{#1}}}
\newcommand{\good}[1]{\textcolor{positive}{#1}}
\newcommand{\bad}[1]{\textcolor{negative}{#1}}

% Standout frame macro: full-bleed dark background for section transitions.
% Beamer-only (uses \begin{frame}/\setbeamercolor/\usebeamerfont) -- do not
% call from an article-class file (e.g. Notes/<CODE>/*-notes.tex). \muted,
% \key, \good, \bad, and \coursecode above are plain xcolor/gdef and are
% safe to use from either class; verified when Notes/article-class support
% was added.
% Expand: \transitionslide{Your transition title}
\newcommand{\transitionslide}[1]{%
  \begingroup
  \setbeamercolor{background canvas}{bg=primary-blue}
  \begin{frame}[plain]
    \centering
    \vfill
    {\usebeamerfont{title}\color{white}\Large #1\par}
    \vfill
  \end{frame}
  \endgroup
}

\coursecode{CS401}
\renewcommand{\thesection}{5.\arabic{section}}

\title{Assignment 5 --- Solution Key\\\large Addressing Modes, CPU Organization, and Bus Structures}
\author{Auqib Hamid Lone\\[0.3em]
  \emph{Department of Computer Science \& Engineering}, \emph{GCET Ganderbal}}
\date{\today}

\begin{document}
\maketitle

\section{Conceptual}

\textbf{Q1.} Direct addressing encodes the effective address $A$ directly in the instruction's address field, and that field is fixed at \emph{compile time} --- it cannot change while the program runs. But \texttt{A[i]}'s address changes on every loop iteration as \texttt{i} changes. \textbf{Displacement (indexed) addressing} should be used instead: $\text{EA} = [R_{\text{base}}] + [R_{\text{index}}]$, where $R_{\text{index}}$ holds \texttt{i}. Because $\text{EA}$ is recomputed from a \emph{register} on every execution of the instruction, the same instruction correctly accesses a different array element on every iteration --- exactly the property direct addressing lacks.

\textbf{Q2.} Indirect addressing: the instruction's address field $A$ points to a memory location, and \emph{that} memory location holds the actual operand's address, so $\text{EA} = M[A]$. This costs two memory reads: one to fetch the pointer from $M[A]$, one to fetch the operand from $M[\text{EA}]$. Register-indirect addressing: a register itself holds the operand's address, so $\text{EA} = [\text{R2}]$ (for example), costing only one memory read (to fetch the operand) plus one register read (to get the address, which is free/fast since it needs no bus transaction). \textbf{Register-indirect is cheaper} because the pointer lives in a register, not memory, eliminating one of indirect addressing's two memory accesses.

\section{Numerical}

\textbf{Q3.} Displacement addressing: $\text{EA} = [\text{R3}] + d = 2048 + 30 = 2078$. \textbf{Cost:} one register read (to obtain $[\text{R3}]$) plus one memory read (to fetch the operand from $M[2078]$) --- two accesses total, one of which is a (cheap) register read.

\textbf{Q4.} \textbf{3 control steps}, identical to \texttt{R1} $\gets$ \texttt{R2+R3}. The single-bus datapath's control-step count depends only on the \emph{shape} of the operation (a two-operand ALU op needs to stage one operand in \texttt{Y}, compute and stage the result in \texttt{Z}, then write back), not on which specific registers are named:
\[
\begin{aligned}
T_1&: Y \gets [\text{R5}] \\
T_2&: Z \gets [Y] + [\text{R6}] \\
T_3&: R4 \gets [Z]
\end{aligned}
\]

\section{Design}

\textbf{Q5.} Subtraction has the identical three-step shape as addition; only the ALU's selected function changes, at $T_2$:
\[
\begin{aligned}
T_1&: Y \gets [\text{R2}] \qquad \text{(R2 onto bus, latched into Y)} \\
T_2&: Z \gets [Y] - [\text{R3}] \qquad \text{(R3 onto bus} \to \text{ALU; ALU subtracts; result} \to Z\text{)} \\
T_3&: R1 \gets [Z] \qquad \text{(Z onto bus, latched into R1)}
\end{aligned}
\]
\textbf{3 control steps}, same as addition, because the bottleneck forcing three steps is structural (the bus can carry only one operand per cycle, and the ALU needs two, then a write-back) --- it does not depend on which ALU operation is selected, only on the fact that it is a two-operand, one-result operation.

\textbf{Q6.} On the three-bus datapath, both operands are read and the result is written back in a single cycle:
\[
T_1: R7 \gets [\text{R4}] + [\text{R5}], \qquad A = [\text{R4}], \ B = [\text{R5}], \ C = \text{result}.
\]
Bus A carries $[\text{R4}]$ into the ALU's first input; Bus B carries $[\text{R5}]$ into the ALU's second input; Bus C carries the ALU's output back into \texttt{R7} via the destination decoder. This instruction requires \textbf{1 control step}, since Bus A and Bus B are separate physical paths that can each drive a different register's contents in the same cycle, and Bus C is not competing with either for the write-back.

\end{document}
